\section{Die vier Schichten sind keine Rangordnung}

\textbf{[SETZUNG]} Die ganze Serie markiert jede Aussage mit einer von vier
Schichten: \textbf{[K]} Kern (aus $\xi$ hergeleitet), \textbf{[B]} Brücke
(algebraisch bewiesen), \textbf{[S]} Skizze (Plausibilität), \textbf{[SETZUNG]}
deklariert. Dieses Dokument sagt, was diese Markierung bedeutet und was sie
nicht bedeutet.

Sie ist eine \emph{Ordnung}, keine \emph{Rangordnung}. Die Schichten sagen, mit
welchem Recht eine Aussage dasteht -- nicht, wie wertvoll sie ist. Eine
\textbf{[SETZUNG]} ist nicht schlechter als ein \textbf{[K]}; sie ist von
anderer Art. Ohne Setzungen gäbe es keinen Kern, aus dem etwas folgen könnte.

\textbf{[B]} Die Verwechslung von Ordnung und Rangordnung ist die häufigste
Fehldeutung geschichteter Darstellungen. Wer \textbf{[S]} als \enquote{noch
nicht ganz [K]} liest, unterstellt, alles strebe zum Kern hin. Das tut es nicht.
Eine Skizze, die als Skizze markiert ist, ist an ihrem Ort vollständig -- sie
behauptet nicht mehr, als sie zeigt, und das ist ihre Stärke.

\section{Was jede Schicht leistet}

\textbf{[SETZUNG]} Die vier Schichten in ihrer eigenen Sprache:

\textbf{[K] Kern.} Aus den drei Setzungen (A020) folgt die Aussage durch
Rechnung oder geometrische Notwendigkeit. Beispiele: der Nichtschließungssatz
(A050), $Q=2/3$ aus $r=\sqrt2$ (A110), $\Delta w = 1$ (A180). Der Kern ist das,
woran die Theorie steht oder fällt.

\textbf{[B] Brücke.} Die Aussage ist algebraisch bewiesen, benutzt aber eine
Übersetzung oder eine bekannte Struktur, die nicht selbst aus $\xi$ folgt.
Beispiel: die Hilbertraum-Darstellung (A070). Eine Brücke trägt alles, was auf
ihren beiden Seiten steht, und beweist nichts Neues -- sie verbindet.

\textbf{[S] Skizze.} Die Aussage ist plausibel, aber nicht durchgerechnet.
Beispiele: der Landauer-Bezug (A180), die Wegzuordnung der Verschränkung (A160).
Eine Skizze ist ein ehrlicher Platzhalter, keine halbe Wahrheit.

\textbf{[SETZUNG] Deklaration.} Die Aussage wird gesetzt, nicht abgeleitet.
Beispiele: die drei Grundsetzungen (A020), die Wahl der charakteristischen Masse
(A145). Eine Setzung ist der Eintrittspunkt, ohne den ein geschlossenes System
nicht anfangen kann (A130).

\section{Die Selbstanwendung}

\textbf{[SETZUNG]} Die Ordnung gilt auch für sich selbst. Die Grundrelation
$\tilde T\cdot m = 1$ ist eine \textbf{[SETZUNG]}, keine \textbf{[K]}. Sie wird
nicht bewiesen; sie ist der Anfang.

Das ist keine Schwäche, sondern die Bedingung der Möglichkeit des Ganzen. Jede
Theorie, die etwas herleitet, leitet es aus etwas her, das sie nicht herleitet.
Der Unterschied zwischen einer sauberen und einer unsauberen Theorie ist nicht,
ob sie Setzungen hat -- jede hat welche --, sondern ob sie sie \emph{als solche
ausweist}.

\textbf{[B]} Genau das leistet die Schichtmarkierung: sie macht die Setzungen
sichtbar, statt sie als Ergebnisse zu tarnen. Der Fehler, den A130 an der
älteren $\alpha$-Darstellung benennt -- eine Kalibrierung als Herleitung
auszugeben --, ist ein Verstoß gegen genau diese Regel. Die Markierung ist das
Mittel, ihn zu vermeiden.

\section{Warum keine Rangordnung möglich ist}

\textbf{[B]} Man könnte versuchen, die Schichten doch zu ordnen: [K] besser als
[B] besser als [S] besser als [SETZUNG]. Der Versuch scheitert an der
Selbstanwendung. Die wertvollste Aussage der Theorie -- $\tilde T\cdot m = 1$ --
wäre dann die geringste. Das ist absurd, und die Absurdität zeigt, dass die
Ordnung horizontal ist, nicht vertikal.

Die Schichten sind vier Weisen, wahr zu sein, nicht vier Grade von Wahrheit.

\section{Was offen bleibt}

\textbf{Erstens ist die Zuordnung im Einzelfall nicht immer scharf.} Manche
Aussagen liegen zwischen [B] und [S]. Die Serie entscheidet solche Fälle
konservativ -- im Zweifel die schwächere Schicht --, aber die Grenze ist nicht
in jedem Fall zwingend.

\textbf{Zweitens ersetzt die Markierung kein Urteil.} Sie sagt, mit welchem
Recht eine Aussage dasteht; ob dieses Recht im Einzelfall korrekt zuerkannt ist,
muss geprüft werden. Die Prüfskripte leisten das für die rechenbaren Fälle.

\section{Ersetzt}

Dieses Dokument nimmt die methodische Note \enquote{Ordnung ohne Rangordnung /
Order without Ranking} auf und macht ihre Unterscheidung (Gesetz, Norm,
Stipulation) zum durchgehenden Prinzip der Serie. Es steht neben A010, das die
Regeln nennt; A200 begründet sie.

\section{Verwendung von KI-Werkzeugen}

Dieses Dokument ist unter Verwendung KI-gestützter Werkzeuge entstanden. Die
Fragestellungen, die Setzungen und alle inhaltlichen Entscheidungen stammen vom
Autor; die Werkzeuge formulieren aus, rechnen nach und erstellen Prüfskripte.
Die Verantwortung für Inhalt und Fehler liegt vollständig beim Autor. Der
ausführliche Wortlaut steht in A010.
